\documentclass[11pt,a4paper]{article}

% ==============================================================================
% 1. LAYOUT & GEOMETRIE
% ==============================================================================
\usepackage[a4paper, left=2.5cm, right=2.5cm, top=3cm, bottom=3cm, headheight=40pt]{geometry}
\usepackage{parskip}
\usepackage{setspace}
\onehalfspacing
\usepackage{multicol}
\usepackage{enumitem}
\setlist[itemize]{label=-, leftmargin=*, noitemsep}
\setcounter{tocdepth}{2}
\setlength{\emergencystretch}{3em}
\tolerance=2000
\AtBeginEnvironment{align}{\small}
\AtBeginEnvironment{align*}{\small}
\AtBeginEnvironment{equation}{\small}
\AtBeginEnvironment{equation*}{\small}
\setlength{\abovedisplayskip}{4pt}
\setlength{\belowdisplayskip}{4pt}
\newcommand{\mysection}[1]{\section*{#1}\addcontentsline{toc}{section}{#1}}
\newcommand{\mysubsection}[1]{\subsection*{#1}\addcontentsline{toc}{subsection}{#1}}

% ==============================================================================
% 2. SPRACHE, FONTS & TYPOGRAFIE
% ==============================================================================
\usepackage[T1]{fontenc}
\usepackage[utf8]{inputenc}
\usepackage[ngerman]{babel}
\usepackage{csquotes}
\usepackage{microtype}

% ==============================================================================
% 3. FARBEN & GRAFIKEN
% ==============================================================================
\usepackage[dvipsnames, table, svgnames]{xcolor}
\usepackage{graphicx}
\usepackage{float}
\usepackage{tikz}
\usepackage{pgfplots}
\pgfplotsset{compat=1.18}

% ==============================================================================
% 4. MATHEMATIK & BAUPHYSIKALISCHE EINHEITEN
% ==============================================================================
\usepackage{amsmath}
\usepackage{amsthm}
\usepackage{amssymb}
\usepackage{amsfonts}
\usepackage{mathtools}
\usepackage{cancel}

\usepackage[
    locale=DE,
    per-mode=symbol,
    separate-uncertainty=true,
    range-units=single
]{siunitx}

\usepackage{physics}
\AtBeginDocument{\RenewCommandCopy\qty\SI}

\providecommand{\eh}[1]{\,[\unit{#1}]}
\newcommand{\gradC}{\,\unit{\degree C}}

% ==============================================================================
% 5. TABELLEN & FORMELBOXEN
% ==============================================================================
\usepackage{booktabs}
\usepackage{tabularx}
\usepackage[most]{tcolorbox}

\newtcolorbox{formelbox}[1]{
    enhanced,
    colback=white,
    colframe=blue!35!black,
    colbacktitle=blue!5!white,
    coltitle=blue!85!black,
    fonttitle=\bfseries\sffamily,
    fontlower=\footnotesize,
    title={#1},
    boxrule=0.8pt,
    arc=4pt,
    left=6pt, right=6pt, top=4pt, bottom=4pt,
    boxsep=0pt,
    toptitle=4pt, bottomtitle=4pt,
    segmentation style={solid, draw=black!20, line width=0.5pt},
    before skip=8pt,
    after skip=8pt,
    unbreakable
}

\newenvironment{symbols}{
    \renewcommand{\arraystretch}{1.15}
    \tabularx{\linewidth}{@{}l l X@{}}
}{
    \endtabularx
}

% ==============================================================================
% 6. KOPF- UND FUSSZEILEN
% ==============================================================================
\usepackage{fancyhdr}
\pagestyle{fancy}
\fancyhf{}
\lhead{\textbf{\large HCU Hamburg}\vspace{0.1cm}}
\rhead{\textbf{\large Bauphysik}\\Formelsammlung}
\lfoot{Bennet Geschke}
\rfoot{Seite \thepage}
\renewcommand{\headrulewidth}{0.4pt}
\renewcommand{\footrulewidth}{0.4pt}

% ==============================================================================
% 7. HYPERREF
% ==============================================================================
\usepackage{hyperref}
\hypersetup{
    pdftitle={Bauphysik - Formelsammlung},
    pdfauthor={Bennet Geschke},
    pdfsubject={Formelsammlung Bauphysik - HCU Hamburg},
    colorlinks=true,
    linkcolor=blue!70!black,
    filecolor=magenta,
    urlcolor=cyan!70!black,
    citecolor=green!60!black
}

\title{\vspace{-2cm}\textbf{\Huge Formelsammlung Bauphysik}}
\author{Bennet Geschke}
\date{\today}

\begin{document}

\maketitle
\vspace{-1.5cm}

\tableofcontents
\vspace{1cm}

\begin{multicols*}{2}

% ==============================================================================
% THEMA 1: WÄRMETRANSPORT
% ==============================================================================

\mysection{1. Wärmetransport}

\mysubsection{1.1 Grundlagen}

\begin{formelbox}{Wärmeenergie}
    \begin{align*}
        Q & = m \cdot c \cdot \Delta T \\
        c_W & = 4,182 \eh{J/(g \cdot K)}
    \end{align*}
    \tcblower
    \begin{symbols}
        $Q$ & \unit{\joule} & Wärmeenergie \\
        $m$ & \unit{\kilo\gram} & Masse \\
        $c$ & \unit{\joule\per\kilo\gram\per\kelvin} & Spezif. Wärmekapazität \\
        $\Delta T$ & \unit{\kelvin} & Temperaturdifferenz \\
        $c_W$ & \unit{J/(g\cdot K)} & Spezif. Wärmekapazität von Wasser
    \end{symbols}
\end{formelbox}

\begin{formelbox}{Wärmestrom (Allgemein)}
    \begin{align*}
        \dot{Q} & = \dot{m} \cdot c \cdot \Delta T
    \end{align*}
    \tcblower
    \begin{symbols}
        $\dot{Q}$ & \unit{\watt} & Wärmestrom \\
        $\dot{m}$ & \unit{\kilo\gram\per\second} & Massenstrom
    \end{symbols}
\end{formelbox}

\mysubsection{1.2 Fourier'sches Gesetz}

\begin{formelbox}{Stationärer Wärmetransport}
    \begin{align*}
        \dot{Q} & = A \cdot \Delta T \cdot \sum \frac{\lambda}{d} \\
        \dot{Q} & = -\lambda \cdot \frac{A}{d} \cdot \Delta T
    \end{align*}
    \tcblower
    \begin{symbols}
        $\lambda$ & \unit{W/(m\cdot K)} & Stoffkonstante, Wärmeleitfähigkeit \\
        $\Delta T$ & \unit{K} & Temperaturdifferenz $T_k - T_w$, $\Delta T < 0$
    \end{symbols}
\end{formelbox}

\begin{formelbox}{Spezifische Wärmeleitfähigkeit}
    \begin{align*}
        \lambda & = \frac{\dot{Q} \cdot d}{A \cdot |\Delta T|}
    \end{align*}
\end{formelbox}

\begin{formelbox}{Wärmestromdichte}
    \begin{align*}
        \dot{Q} & = q \cdot A \\
        q & = \frac{\lambda}{d} \cdot \Delta T
    \end{align*}
    \tcblower
    \begin{symbols}
        $q$ & \unit{W/m^2} & Wärmestromdichte
    \end{symbols}
\end{formelbox}

\mysubsection{1.3 Wärmedurchgang und Wärmedurchlass}

\begin{formelbox}{Wärmedurchlasswiderstand einer Schicht}
    \begin{align*}
        R & = \frac{d}{\lambda} \eh{m^2 K/W}
    \end{align*}
    \tcblower
    \begin{symbols}
        $R$ & \unit{m^2 K/W} & Wärmedurchlasswiderstand einer Schicht
    \end{symbols}
\end{formelbox}

\begin{formelbox}{Wärmedurchlasskoeffizient einer Schicht}
    \begin{align*}
        \Lambda & = \frac{\lambda}{d} \eh{W/(m^2 K)}
    \end{align*}
    \tcblower
    \begin{symbols}
        $\Lambda$ & \unit{W/(m^2 K)} & Wärmedurchlasskoeffizient einer Schicht
    \end{symbols}
\end{formelbox}

\begin{formelbox}{Wärmedurchgangswiderstand gesamt}
    \begin{align*}
        R_T & = R_{si} + \frac{d_1}{\lambda_1} + \dots + \frac{d_n}{\lambda_n} + R_{se} \\
        R_T & = R_{si} + \sum R_i + R_{se}
    \end{align*}
    \tcblower
    \begin{symbols}
        $R_T$ & \unit{m^2 K/W} & Wärmedurchgangswiderstand gesamt \\
        $R_{si}$ & \unit{m^2 K/W} & Wärmeübergangswiderstand innen ($0,1$ $\uparrow$, $0,13$ $\rightarrow$, $0,17$ $\downarrow$) \\
        $R_{se}$ & \unit{m^2 K/W} & Wärmeübergangswiderstand außen ($0,04$) \\
        $R_i$ & \unit{m^2 K/W} & Wärmedurchlasswiderstand der Einzelschicht
    \end{symbols}
\end{formelbox}

\begin{formelbox}{Wärmedurchgangskoeffizient U}
    \begin{align*}
        U & = \frac{1}{R_T} \eh{W/(m^2 K)} \\
        U & = \frac{q}{\vartheta_{ai} - \vartheta_{ae}} \eh{W/(m^2 K)}
    \end{align*}
    \tcblower
    \begin{symbols}
        $U$ & \unit{W/(m^2 K)} & Wärmedurchgangskoeffizient \\
        $\vartheta_{ai}$ & \unit{K} & Innenlufttemperatur \\
        $\vartheta_{ae}$ & \unit{K} & Außenlufttemperatur
    \end{symbols}
\end{formelbox}

\begin{formelbox}{Wärmestrom durch Wand}
    \begin{align*}
        \dot{Q} & = U \cdot A \cdot (\vartheta_{ai} - \vartheta_{ae}) \eh{W}
    \end{align*}
\end{formelbox}

\begin{formelbox}{Transmissionswärmeverlust}
    \begin{align*}
        H_T & = U \cdot A
    \end{align*}
    \tcblower
    \begin{symbols}
        $H_T$ & \unit{W/K} & Transmissionswärmeverlust-Koeffizient
    \end{symbols}
\end{formelbox}

\mysubsection{1.4 Temperaturverlauf im Bauteil}

\begin{formelbox}{Temperaturdifferenz einer Teilschicht}
    \begin{align*}
        \frac{\Delta T_i}{\Delta T_{ges}} & = \frac{R_i}{R_T}
    \end{align*}
    \tcblower
    \begin{symbols}
        $\Delta T_i$ & \unit{K} & Temperaturdifferenz der Einzelschicht \\
        $\Delta T_{ges}$ & \unit{K} & Temperaturdifferenz der Gesamtschicht
    \end{symbols}
\end{formelbox}

\begin{formelbox}{Temperaturverlauf}
    \begin{align*}
        \theta_k & = \theta_i - \sum \Delta T_i \\
        \theta_k & = \theta_e + \sum \Delta T_i
    \end{align*}
    \tcblower
    \begin{symbols}
        $\theta_k$ & \unit{K} & Temperatur an der Schichtgrenze \\
        $\theta_i$ & \unit{K} & Innentemperatur \\
        $\theta_e$ & \unit{K} & Außentemperatur
    \end{symbols}
\end{formelbox}

\begin{formelbox}{Temperaturfaktor (Schimmelkriterium)}
    \begin{align*}
        f_{Rsi} & = \frac{\vartheta_{si} - \vartheta_{ae}}{\vartheta_{ai} - \vartheta_{ae}} > 0,7
    \end{align*}
    \tcblower
    \begin{symbols}
        $f_{Rsi}$ & --- & Temperaturfaktor an der Innenoberfläche \\
        $\vartheta_{si}$ & \unit{K} & Innenoberflächentemperatur
    \end{symbols}
\end{formelbox}

\mysubsection{1.5 Mittelwerte und Wärmebrücken}

\begin{formelbox}{Wärmedurchgangswiderstand (mittel)}
    \begin{align*}
        R_T & = \frac{R_{T,min} + R_{T,max}}{2} \\
        U_m & = \frac{1}{R_T}
    \end{align*}
    \tcblower
    \begin{symbols}
        $R_{T,min}$ & \unit{m^2 K/W} & Minimaler Wärmedurchgangswiderstand \\
        $R_{T,max}$ & \unit{m^2 K/W} & Maximaler Wärmedurchgangswiderstand \\
        $U_m$ & \unit{W/(m^2 K)} & Mittlerer Wärmedurchgangskoeffizient
    \end{symbols}
\end{formelbox}

\begin{formelbox}{Wärmedurchgang (Nebeneinanderliegende Flächen)}
    \begin{align*}
        R_{T,max} & = \frac{1}{U_{min}} \\
        U_{min} & = \frac{\sum U_j \cdot A_j}{\sum A_j} \\
        U_j & = \frac{1}{R_j} \\
        R_j & = R_{si} + \sum R_i + R_{se}
    \end{align*}
    \tcblower
    \begin{symbols}
        $U_j$ & \unit{W/(m^2 K)} & Einzelne Wärmedurchgangskoeffizienten \\
        $A_j$ & \unit{m^2} & Einzelne Flächen
    \end{symbols}
\end{formelbox}

\begin{formelbox}{Wärmedurchgang (Reihenschaltung)}
    \begin{align*}
        R_{T,min} & = R_{si} + \sum \frac{1}{\Lambda_i} + R_{se} \\
        \Lambda_i & = \frac{\sum \Lambda_j \cdot A_j}{\sum A_j}
    \end{align*}
\end{formelbox}

\begin{formelbox}{Wärmebrückenzuschlag}
    \begin{align*}
        U_{ges} & = U + \Delta U_{WB} \\
        \Delta U_{WB} & = \frac{1}{A_{ges}} \cdot \left( \sum (l_j \cdot \psi_j) + \sum \chi_k \right)
    \end{align*}
    \tcblower
    \begin{symbols}
        $\Delta U_{WB}$ & \unit{W/(m^2 K)} & Wärmebrückenzuschlag \\
        & & $= 0,10$ (ohne Nachweis), $0,05$ (DIN 4108-2) \\
        $l_j$ & \unit{m} & Länge der Wärmebrücke \\
        $\psi_j$ & \unit{W/(m K)} & Psi-Wert der Wärmebrücke \\
        $\chi_k$ & \unit{W/K} & Chi-Wert der Wärmebrücke
    \end{symbols}
\end{formelbox}

% ==============================================================================
% THEMA 2: WÄRMESTRAHLUNG
% ==============================================================================

\mysection{2. Wärmestrahlung}

\mysubsection{2.1 Grundlagen der Strahlung}

\begin{formelbox}{Stefan-Boltzmann-Gesetz (schwarzer Körper)}
    \begin{align*}
        q & = \sigma \cdot T^4 \\
        \sigma & = 5,67 \cdot 10^{-8} \eh{W/(m^2 K^4)}
    \end{align*}
    \tcblower
    \begin{symbols}
        $q$ & \unit{W/m^2} & Abgestrahlte Intensität \\
        $\sigma$ & \unit{W/(m^2 K^4)} & Stefan-Boltzmann-Konstante
    \end{symbols}
\end{formelbox}

\begin{formelbox}{Stefan-Boltzmann-Gesetz (nicht-schwarzer Körper)}
    \begin{align*}
        q & = \sigma \cdot \epsilon \cdot T^4
    \end{align*}
    \tcblower
    \begin{symbols}
        $\epsilon$ & --- & Emissionsgrad ($\epsilon < 1$ für nicht-schwarze Körper)
    \end{symbols}
\end{formelbox}

\begin{formelbox}{Wien'sches Verschiebungsgesetz}
    \begin{align*}
        \lambda_{max} & = \frac{2898 \,\mu\text{m} \cdot K}{T}
    \end{align*}
    \tcblower
    \begin{symbols}
        $\lambda_{max}$ & \unit{\mu m} & Wellenlänge des Intensitätsmaximums
    \end{symbols}
\end{formelbox}

\begin{formelbox}{Absorptions-, Emissions-, Reflexions- \& Transmissionsgrad}
    \begin{align*}
        \alpha & = \frac{I_A}{I}, \quad \rho = \frac{I_R}{I}, \quad \tau = \frac{I_T}{I} \\
        \alpha & = \epsilon \quad \text{(Kirchhoff'sches Gesetz)}
    \end{align*}
    \tcblower
    \begin{symbols}
        $\alpha$ & --- & Absorptionsgrad \\
        $\rho$ & --- & Reflexionsgrad \\
        $\tau$ & --- & Transmissionsgrad \\
        $\epsilon$ & --- & Emissionsgrad
    \end{symbols}
\end{formelbox}

\mysubsection{2.2 Strahlungsaustausch}

\begin{formelbox}{Strahlungsaustausch planparalleler Flächen}
    \begin{align*}
        q & = \sigma \cdot \epsilon_{pp} \cdot (T_1^4 - T_2^4) \\
        \epsilon_{pp} & = \frac{1}{\frac{1}{\epsilon_1} + \frac{1}{\epsilon_2} - 1}
    \end{align*}
    \tcblower
    \begin{symbols}
        $q$ & \unit{W/m^2} & Austausch-Wärmestromdichte \\
        $\epsilon_{pp}$ & --- & Gemeinsamer Emissionsgrad (gilt für unendlich ausgedehnte planparallele Oberflächen)
    \end{symbols}
\end{formelbox}

\begin{formelbox}{Wärmestrahlung Einzelfläche}
    \begin{align*}
        q_{10} & = \sigma \cdot \epsilon_i \cdot T_i^4 \\
        \epsilon_i & = 1 - \rho_i
    \end{align*}
\end{formelbox}

\begin{formelbox}{Solarkonstante}
    \begin{align*}
        q_{max} & = 1,3 \eh{kW/m^2} \quad \text{(oberhalb Atmosphäre)} \\
        q_{max} & = 1,0 \eh{kW/m^2} \quad \text{(unterhalb Atmosphäre)}
    \end{align*}
\end{formelbox}

% ==============================================================================
% THEMA 3: SOLARE GEWINNE UND WÄRMESCHUTZ
% ==============================================================================

\mysection{3. Solarer Wärmeschutz}

\mysubsection{3.1 Kenngrößen für Fenster}

\begin{formelbox}{Gesamtenergiedurchlassgrad g}
    Typische Werte (DIN 4108-4):
    \begin{align*}
        \text{Einfachverglasung} & : g = 0,87 \\
        \text{Doppelverglasung} & : g = 0,78 \\
        \text{Wärmeschutzglas (selektiv)} & : g = 0,50 - 0,70 \\
        \text{Dreifachverglasung normal} & : g = 0,60 - 0,70 \\
        \text{Dreifachverglasung (2x selektiv)} & : g = 0,35 - 0,50 \\
        \text{Sonnenschutzglas} & : g = 0,20 - 0,50
    \end{align*}
\end{formelbox}

\begin{formelbox}{Lichttransmissionsgrad}
    \begin{align*}
        \tau_r & \text{(Richtwerte)} \\
        \text{Einfachglas} & : \tau_r = 0,90 \\
        \text{Zweifachglas ohne Beschichtung} & : \tau_r = 0,82 \\
        \text{Dreifachglas ohne Beschichtung} & : \tau_r = 0,75 \\
        \text{Wärmedämmglas 2-fach} & : \tau_r = 0,60 - 0,78
    \end{align*}
\end{formelbox}

\begin{formelbox}{Abminderungsfaktor Sonnenschutz}
    \begin{align*}
        F_c & = \frac{g_{tot}}{g}
    \end{align*}
    \tcblower
    \begin{symbols}
        $g_{tot}$ & --- & Tatsächlicher Gesamtenergiedurchlassgrad mit Sonnenschutz \\
        $g$ & --- & Gesamtenergiedurchlassgrad des Glases ohne Sonnenschutz
    \end{symbols}
\end{formelbox}

\begin{formelbox}{Sonnenschutz-Faktoren (nach Tabelle)}
    \begin{align*}
        \text{Innenliegend / zwischen Scheiben:} \\
        \text{weiß/hochreflektierend} &: F_c = 0,65-0,70 \\
        \text{helle Farben} &: F_c = 0,75-0,80 \\
        \text{dunkle Farben} &: F_c = 0,85-0,90 \\
        \text{Außenliegend:} \\
        \text{Fensterläden geschlossen} &: F_c = 0,10-0,15 \\
        \text{Jalousie 45°} &: F_c = 0,25-0,30 \\
        \text{Markise} &: F_c = 0,25-0,30
    \end{align*}
\end{formelbox}

\mysubsection{3.2 Sonneneintragskennwert}

\begin{formelbox}{Vorhandener Sonneneintragskennwert S}
    \begin{align*}
        S_{vorh} & = \frac{\sum_{j=1}^{N} A_{w,j} \cdot g_{i,j} \cdot F_{C,j}}{A_G}
    \end{align*}
    \tcblower
    \begin{symbols}
        $A_w$ & \unit{m^2} & Fensterfläche (lichte Rohbaumaße) \\
        $g$ & --- & Gesamtenergiedurchlassgrad \\
        $F_C$ & --- & Abminderungsfaktor Sonnenschutz \\
        $A_G$ & \unit{m^2} & Grundfläche (lichte Rohbaumaße) \\
        $N$ & --- & Anzahl der Seitenwände (meist 4)
    \end{symbols}
\end{formelbox}

\begin{formelbox}{Zulässiger Sonneneintragskennwert S}
    \begin{align*}
        S_{zul} & = \sum_{i=1}^{6} S_i
    \end{align*}
    \tcblower
    \begin{symbols}
        $S_1$ & --- & Nachtlüftung und Bauart \\
        $S_2$ & --- & Fensteranteil \\
        $S_3$ & --- & Sonnenschutzglas \\
        $S_4$ & --- & Fensterneigung \\
        $S_5$ & --- & Orientierung \\
        $S_6$ & --- & Einsatz passiver Kühlung
    \end{symbols}
\end{formelbox}

\begin{formelbox}{Fensteranteil $S_2$}
    \begin{align*}
        S_2 & = a - b \cdot f_{WG} \\
        f_{WG} & = \frac{A_W}{A_G}
    \end{align*}
    \tcblower
    \begin{symbols}
        $a$ & --- & Wohngebäude: $0,060$; Nichtwohngebäude: $0,030$ \\
        $b$ & --- & Wohngebäude: $0,231$; Nichtwohngebäude: $0,115$
    \end{symbols}
\end{formelbox}

\begin{formelbox}{Bagatellbedingung}
    \begin{align*}
        f_{WG} & \leq \text{erf. } f_{WG}
    \end{align*}
\end{formelbox}

\begin{formelbox}{Effektive Wärmespeicherkapazität}
    \begin{align*}
        C_{wirk} & = \sum_i m_i \cdot c_i = \sum_i \rho_i \cdot d_i \cdot A_i \cdot c_i
    \end{align*}
    \tcblower
    \begin{symbols}
        $C_{wirk}$ & \unit{Wh/K} & Effektive Wärmespeicherkapazität \\
        $\rho_i$ & \unit{kg/m^3} & Rohdichte der Schicht \\
        $d_i$ & \unit{m} & Schichtdicke \\
        $A_i$ & \unit{m^2} & Fläche der Schicht \\
        $c_i$ & \unit{Wh/(kg\cdot K)} & Spezifische Wärmekapazität
    \end{symbols}
    \begin{itemize}
        \item Mittlere Bauart: $50 \leq C_{wirk}/A_G \leq 130$ \unit{Wh/(K\cdot m^2)}
    \end{itemize}
\end{formelbox}

% ==============================================================================
% THEMA 4: FEUCHTE UND TAUPUNKT
% ==============================================================================

\mysection{4. Feuchte und Taupunkt}

\mysubsection{4.1 Zustandsgleichungen}

\begin{formelbox}{Zustandsgleichung für ideale Gase}
    \begin{align*}
        pV & = nRT \\
        pV & = m R_S T \\
        p & = \rho R_S T
    \end{align*}
    \tcblower
    \begin{symbols}
        $p$ & \unit{Pa} & Druck \\
        $V$ & \unit{m^3} & Volumen \\
        $T$ & \unit{K} & Temperatur \\
        $R$ & \unit{J/(mol\cdot K)} & Allgemeine Gaskonstante ($8,31$) \\
        $R_S$ & \unit{J/(kg\cdot K)} & Spezifische Gaskonstante \\
        $R_D$ & \unit{J/(kg\cdot K)} & Gaskonstante für Wasserdampf ($461,5$)
    \end{symbols}
\end{formelbox}

\mysubsection{4.2 Sättigungsdampfdruck}

\begin{formelbox}{Sättigungsdampfdruck $p_s$ (DIN 4108-3)}
    \begin{align*}
        p_s & = 610,5 \cdot \exp\left(\frac{17,269 \cdot \vartheta}{237,3 + \vartheta}\right) \quad \text{für } \vartheta \geq 0\gradC \\
        p_s & = 610,5 \cdot \exp\left(\frac{21,875 \cdot \vartheta}{265,5 + \vartheta}\right) \quad \text{für } \vartheta < 0\gradC
    \end{align*}
\end{formelbox}

\mysubsection{4.3 Luftfeuchte}

\begin{formelbox}{Absolute Luftfeuchte}
    \begin{align*}
        \rho_D & \text{ in } \unit{kg/m^3} \quad \text{(in Luft enthaltenes Wasser)}
    \end{align*}
\end{formelbox}

\begin{formelbox}{Relative Luftfeuchte}
    \begin{align*}
        \Phi & = \frac{\rho_D}{\rho_s} = \frac{p_D}{p_s}
    \end{align*}
    \tcblower
    \begin{symbols}
        $\Phi$ & --- & Relative Luftfeuchte \\
        $\rho_D$ & \unit{kg/m^3} & Dampfdichte \\
        $\rho_s$ & \unit{kg/m^3} & Sättigungsdampfdichte \\
        $p_D$ & \unit{Pa} & Dampfdruck \\
        $p_s$ & \unit{Pa} & Sättigungsdampfdruck
    \end{symbols}
\end{formelbox}

\begin{formelbox}{Dampfdruck aus Dampfdichte}
    \begin{align*}
        p_D & = \rho_D \cdot R_D \cdot T
    \end{align*}
\end{formelbox}

\mysubsection{4.4 Taupunktberechnung}

\begin{formelbox}{Taupunkt $\tau$ aus $\vartheta_1$ und $\Phi_1$}
    \begin{align*}
        \tau & = \frac{237,3 \cdot \left(\frac{17,269 \cdot \vartheta_1}{237,3 + \vartheta_1} + \ln \Phi_1\right)}{17,269 - \left(\frac{17,269 \cdot \vartheta_1}{237,3 + \vartheta_1} + \ln \Phi_1\right)}
    \end{align*}
\end{formelbox}

\begin{formelbox}{Temperatur $\vartheta_2$ aus $\vartheta_1$, $\Phi_1$ und $\Phi_2$}
    \begin{align*}
        \vartheta_2 & = \frac{237,3 \cdot \left(\frac{17,269 \cdot \vartheta_1}{237,3 + \vartheta_1} + \ln\left(\frac{\Phi_1}{\Phi_2}\right)\right)}{17,269 - \left(\frac{17,269 \cdot \vartheta_1}{237,3 + \vartheta_1} + \ln\left(\frac{\Phi_1}{\Phi_2}\right)\right)}
    \end{align*}
\end{formelbox}

\begin{formelbox}{Relative Feuchte $\Phi_2$ aus $\vartheta_1$, $\vartheta_2$ und $\Phi_1$}
    \begin{align*}
        \Phi_2 & = \Phi_1 \cdot \exp\left(\frac{17,269 \cdot \vartheta_1}{237,3 + \vartheta_1} - \frac{17,269 \cdot \vartheta_2}{237,3 + \vartheta_2}\right)
    \end{align*}
\end{formelbox}

\begin{formelbox}{Allgemeine Feuchtebeziehung}
    \begin{align*}
        \Phi_1 \cdot p_{s1}(\theta_1) & = \Phi_2 \cdot p_{s2}(\theta_2)
    \end{align*}
\end{formelbox}

\mysubsection{4.5 Feuchtegleichgewicht und Dampfdiffusion}

\begin{formelbox}{Feuchtegleichgewicht}
    \begin{align*}
        \Phi_i \cdot p_{si} - \Phi_a \cdot p_{sa} & = \frac{\dot{m}_D \cdot R_D \cdot T_i}{N_L \cdot V}
    \end{align*}
\end{formelbox}

\begin{formelbox}{Dampfstrom}
    \begin{align*}
        \dot{m}_D & = \delta_0 \cdot \frac{\Delta p}{s_{d\tau}} \\
        s_{d\tau} & = \sum \mu_i \cdot d_i \\
        \Delta p & = \Phi_i \cdot p_{si} - \Phi_e \cdot p_{se}
    \end{align*}
    \tcblower
    \begin{symbols}
        $\dot{m}_D$ & \unit{kg/s} & Dampfstrom \\
        $\delta_0$ & --- & Dampfdiffusionskoeffizient \\
        $s_{d\tau}$ & \unit{m} & Diffusionsäquivalente Luftschichtdicke \\
        $\mu_i$ & --- & Diffusionswiderstandszahl \\
        $d_i$ & \unit{m} & Schichtdicke
    \end{symbols}
\end{formelbox}

\begin{formelbox}{Raumfeuchte}
    \begin{align*}
        \Phi & = \frac{p_d}{p_s} \\
        p_d & = \frac{m_D \cdot R_D \cdot T}{V}
    \end{align*}
\end{formelbox}

% ==============================================================================
% THEMA 5: LÜFTUNGSWÄRMEVERLUSTE
% ==============================================================================

\mysection{5. Lüftungswärmeverluste}

\begin{formelbox}{Lüftungswärmeverluste (DIN 4108-6)}
    \begin{align*}
        H_V & = 0,19 \cdot V_e \quad \text{(ohne Dichtheitsprüfung)} \\
        H_V & = 0,163 \cdot V_e \quad \text{(mit Dichtheitsprüfung)}
    \end{align*}
    \tcblower
    \begin{symbols}
        $H_V$ & \unit{W/K} & Lüftungswärmeverlustkoeffizient \\
        $V_e$ & \unit{m^3} & Gebäudevolumen
    \end{symbols}
\end{formelbox}

\begin{formelbox}{Luftdichtheit nach DIN 4108-6}
    \begin{align*}
        \text{Einfamilienwohnhaus} &: n_{50} = 1,0 - 20,0 \quad (1/h) \\
        \text{Mehrfamilienwohnhaus} &: n_{50} = 0,5 - 10,0 \quad (1/h)
    \end{align*}
    \begin{itemize}
        \item sehr dicht: $0,5 - 2,0$ (MFH), $1,0 - 3,0$ (EFH)
        \item mittel dicht: $2,0 - 4,0$ (MFH), $3,0 - 8,0$ (EFH)
        \item wenig dicht: $4,0 - 10,0$ (MFH), $8,0 - 20,0$ (EFH)
    \end{itemize}
\end{formelbox}

% ==============================================================================
% THEMA 6: SCHALL
% ==============================================================================

\mysection{6. Schall}

\mysubsection{6.1 Pegelgrößen}

\begin{formelbox}{Allgemeine Pegelgrößen}
    \begin{align*}
        L_w & = 10 \cdot \log\left(\frac{P}{P_0}\right), \quad P_0 = 10^{-12} \unit{W} \\
        L_I & = 10 \cdot \log\left(\frac{I}{I_0}\right), \quad I_0 = 10^{-12} \unit{W/m^2} \\
        L_p & = 10 \cdot \log\left(\frac{p^2}{p_0^2}\right), \quad p_0 = 2 \cdot 10^{-5} \unit{Pa} \\
        L_v & = 10 \cdot \log\left(\frac{v^2}{v_0^2}\right), \quad v_0 = 5 \cdot 10^{-8} \unit{m/s}
    \end{align*}
\end{formelbox}

\begin{formelbox}{Schallpegeladdition}
    \begin{align*}
        L_{ges} & = 10 \cdot \log\left(\sum_{i=1}^{n} 10^{0,1 \cdot L_i}\right)
    \end{align*}
\end{formelbox}

\begin{formelbox}{A-bewerteter Schallpegel}
    \begin{align*}
        L_{dB-A} & = L_{dB} + \Delta L
    \end{align*}
    \begin{itemize}
        \item A-Bewertung nach DIN EN 61672-1
        \item $\Delta L$ frequenzabhängig nach Tabellen
    \end{itemize}
\end{formelbox}

\mysubsection{6.2 Schallausbreitung}

\begin{formelbox}{Punktquelle (Kugelquelle)}
    \begin{align*}
        I(r) & = \frac{P}{4\pi r^2} \\
        L(r) & = L_w - 10 \cdot \log(4\pi r^2) \\
        L(1\unit{m}) & = L_w - 11 \unit{dB} \\
        L(r_2) & = L(r_1) - 20 \cdot \log\left(\frac{r_2}{r_1}\right)
    \end{align*}
    \tcblower
    \begin{symbols}
        $I$ & \unit{W/m^2} & Schallintensität \\
        $P$ & \unit{W} & Schallleistung \\
        $r$ & \unit{m} & Abstand \\
        $\Delta L$ & \unit{dB} & Pegelabnahme: $20\log(r_2/r_1)$
    \end{symbols}
\end{formelbox}

\begin{formelbox}{Linienquelle}
    \begin{align*}
        I & = \frac{P'}{4a} \\
        L(1\unit{m}) & = L_w - 11 \unit{dB} \\
        L(r_2) & = L(r_1) - 10 \cdot \log\left(\frac{r_2}{r_1}\right)
    \end{align*}
    \tcblower
    \begin{symbols}
        $P'$ & \unit{W/m} & Schallleistung pro Meter \\
        $a$ & \unit{m} & Senkrechter Abstand zur Linienquelle
    \end{symbols}
\end{formelbox}

\begin{formelbox}{Allgemeine Pegeländerung}
    \begin{align*}
        L(r_2) & = L(r_1) - K \cdot 10 \cdot \log\left(\frac{r_2}{r_1}\right) \\
        K & = 1 \text{ für Linienquelle, } K = 2 \text{ für Punktquelle}
    \end{align*}
\end{formelbox}

\mysubsection{6.3 Schallabsorption in Luft}

\begin{formelbox}{Luftabsorptionskoeffizient}
    \begin{align*}
        m & \approx \frac{0,04}{\Phi} \cdot f^{1,6} \eh{1/m} \\
        \Delta L & = -4,343 \cdot m \cdot r
    \end{align*}
    \tcblower
    \begin{symbols}
        $m$ & \unit{1/m} & Luftabsorptionskoeffizient \\
        $\Phi$ & \% & Relative Luftfeuchte \\
        $f$ & \unit{kHz} & Frequenz
    \end{symbols}
\end{formelbox}

\begin{formelbox}{Pegelminderung durch Luftabsorption}
    \begin{align*}
        L(r_2) & = L(r_1) - 20 \cdot \log\left(\frac{r_2}{r_1}\right) - 4,343 \cdot m \cdot (r_2 - r_1)
    \end{align*}
\end{formelbox}

\mysubsection{6.4 Reflexion}

\begin{formelbox}{Reflexions- und Transmissionsfaktor}
    \begin{align*}
        r & = \frac{Z_2 - Z_1}{Z_2 + Z_1} \\
        t & = 1 + r = \frac{2Z_2}{Z_2 + Z_1}
    \end{align*}
    \tcblower
    \begin{symbols}
        $Z$ & \unit{kg/(m^2\cdot s)} & Schallkennimpedanz
    \end{symbols}
\end{formelbox}

\begin{formelbox}{Reflexionsfaktor bei schrägem Einfall}
    \begin{align*}
        r & = \frac{Z_2 \cos\phi - Z_1}{Z_2 \cos\phi + Z_1}
    \end{align*}
\end{formelbox}

\begin{formelbox}{Gesamtschallpegel bei Reflexion}
    \begin{align*}
        L_{dir} & = L_0 - K \cdot 10 \cdot \log\left(\frac{r}{r_0}\right) \\
        L_{refl} & = L_0 - K \cdot 10 \cdot \log\left(\frac{r}{r_0}\right) + 10 \cdot \log(1 - \alpha) \\
        L_{\Sigma} & = 10 \cdot \log\left(10^{0,1 L_{dir}} + 10^{0,1 L_{refl}}\right)
    \end{align*}
    \tcblower
    \begin{symbols}
        $L_0$ & \unit{dB} & Pegel im Abstand $r_0$ \\
        $\alpha$ & --- & Absorptionsgrad \\
        $\rho = 1 - \alpha$ & --- & Reflexionsgrad
    \end{symbols}
\end{formelbox}

\mysubsection{6.5 Schallbeugung}

\begin{formelbox}{Umwegregel}
    \begin{align*}
        z & = A + B - (a + b) \\
        z & = \frac{h^2}{2} \cdot \left(\frac{1}{a} + \frac{1}{b}\right)
    \end{align*}
    \tcblower
    \begin{symbols}
        $z$ & \unit{m} & Umweg \\
        $A+B$ & \unit{m} & Abstand über Hindernis \\
        $a+b$ & \unit{m} & Direkter Abstand \\
        $h$ & \unit{m} & Höhe des Hindernisses
    \end{symbols}
\end{formelbox}

\begin{formelbox}{Fresnelzahl}
    \begin{align*}
        N & = \frac{z}{\lambda/2} = \frac{2z}{\lambda}
    \end{align*}
    \tcblower
    \begin{symbols}
        $N$ & --- & Fresnelzahl \\
        $\lambda$ & \unit{m} & Wellenlänge
    \end{symbols}
\end{formelbox}

\begin{formelbox}{Pegelminderung durch Umweg}
    \begin{align*}
        T & = \frac{1}{3(1 + 2\pi N)} \\
        \Delta L_z & = -10 \cdot \log(3 + 20N)
    \end{align*}
    \tcblower
    \begin{symbols}
        $T$ & --- & Transmissionsgrad \\
        $\Delta L_z$ & \unit{dB} & Pegelminderung durch Umweg
    \end{symbols}
\end{formelbox}

\begin{formelbox}{Gemittelte Schallbeugung}
    \begin{align*}
        \text{Straßenlärm, Linienquelle:} & \quad \Delta L_z = 10 \cdot \log\left(1 + 80 \cdot \frac{z}{\unit{m}}\right) \\
        \text{Schienenlärm, Punktquelle:} & \quad \Delta L_z = 10 \cdot \log\left(1 + 80 \cdot \frac{z}{\unit{m}}\right)
    \end{align*}
\end{formelbox}

\mysubsection{6.6 Raumakustik}

\begin{formelbox}{Äquivalente Absorptionsfläche}
    \begin{align*}
        A_f & = \sum_{i=1}^{n} S_i \cdot \alpha_{i,f} + 4m_f \cdot V
    \end{align*}
    \tcblower
    \begin{symbols}
        $A_f$ & \unit{m^2} & Äquivalente Absorptionsfläche \\
        $S_i$ & \unit{m^2} & Oberfläche des Bauteils \\
        $\alpha_{i,f}$ & --- & Absorptionsgrad bei Frequenz $f$ \\
        $m_f$ & \unit{1/m} & Luftabsorptionskoeffizient \\
        $V$ & \unit{m^3} & Raumvolumen
    \end{symbols}
\end{formelbox}

\begin{formelbox}{Mittlerer Absorptionsgrad}
    \begin{align*}
        \overline{\alpha} & = \frac{1}{S} \left(\sum_{i=1}^{n} S_i \cdot \alpha_{i,f} + 4m_f \cdot V\right)
    \end{align*}
    \tcblower
    \begin{symbols}
        $S$ & \unit{m^2} & Gesamtoberfläche
    \end{symbols}
\end{formelbox}

\begin{formelbox}{Nachhallzeit nach Sabine}
    \begin{align*}
        T_{60,Sab} & = 0,161 \cdot \frac{V}{A_f}
    \end{align*}
\end{formelbox}

\begin{formelbox}{Nachhallzeit nach Eyring}
    \begin{align*}
        T_{60,Eyr} & = 0,161 \cdot \frac{V}{S \cdot (-\ln(1 - \overline{\alpha}))}
    \end{align*}
\end{formelbox}

\begin{formelbox}{Zusammenhang Sabine/Eyring}
    \begin{align*}
        \overline{\alpha}_{Eyring} & = -\ln(1 - \overline{\alpha}_{Sabine})
    \end{align*}
\end{formelbox}

\begin{formelbox}{Hallradius}
    \begin{align*}
        r_h & \text{ ist der Abstand, bei dem Direktschall und Diffusschall gleich laut sind}
    \end{align*}
\end{formelbox}

\begin{formelbox}{Gesamtintensität im semi-diffusen Schallfeld}
    \begin{align*}
        I & = \frac{P}{4\pi r^2} + \frac{4P \cdot (1 - \overline{\alpha})}{A}
    \end{align*}
    \tcblower
    \begin{symbols}
        $I$ & \unit{W/m^2} & Gesamtintensität \\
        $P$ & \unit{W} & Schallleistung
    \end{symbols}
\end{formelbox}

\begin{formelbox}{Gesamtpegel im semi-diffusen Schallfeld}
    \begin{align*}
        L & = L_w + 10 \cdot \log\left(\frac{1}{4\pi r^2} + \frac{4(1 - \overline{\alpha})}{A}\right)
    \end{align*}
\end{formelbox}

\begin{formelbox}{Intensität im stationären diffusen Schallfeld}
    \begin{align*}
        I & = \frac{4P}{A}
    \end{align*}
\end{formelbox}

\begin{formelbox}{Pegelkorrektur bei Absorption}
    \begin{align*}
        \Delta L & = -10 \cdot \log\left(\frac{A_{nach}}{A_{vor}}\right)
    \end{align*}
\end{formelbox}

\begin{formelbox}{Pegel im diffusen Schallfeld}
    \begin{align*}
        L & = L_w + 6 - 10 \cdot \log\left(\frac{A}{1\unit{m^2}}\right)
    \end{align*}
\end{formelbox}

\mysubsection{6.7 Schalldämmung}

\begin{formelbox}{Schalldämmmaß $R$}
    \begin{align*}
        R & = -10 \cdot \log(T) \\
        R & = L_S - L_E + 10 \cdot \log\left(\frac{S}{A_E}\right)
    \end{align*}
    \tcblower
    \begin{symbols}
        $R$ & \unit{dB} & Schalldämmmaß \\
        $T$ & --- & Transmissionsgrad \\
        $L_S$ & \unit{dB} & Pegel im Senderraum \\
        $L_E$ & \unit{dB} & Pegel im Empfangsraum \\
        $S$ & \unit{m^2} & Fläche der Trennwand \\
        $A_E$ & \unit{m^2} & Äquivalente Absorptionsfläche im Empfangsraum
    \end{symbols}
\end{formelbox}

\begin{formelbox}{Massegesetz (einschalige, biegeweiche Wand)}
    \begin{align*}
        R & = 20 \cdot \log\left(\frac{m'}{\unit{kg/m^2}} \cdot \frac{f}{\unit{Hz}}\right) - 45 \unit{dB}
    \end{align*}
    \tcblower
    \begin{symbols}
        $m'$ & \unit{kg/m^2} & Flächenbezogene Masse \\
        $f$ & \unit{Hz} & Frequenz
    \end{symbols}
\end{formelbox}

\begin{formelbox}{Schätzformel bei 700 Hz Mittenfrequenz}
    \begin{align*}
        R & = 12 + 20 \cdot \log\left(\frac{m'}{\unit{kg/m^2}}\right)
    \end{align*}
\end{formelbox}

\begin{formelbox}{Flächenbezogene Masse}
    \begin{align*}
        m' & = \rho \cdot d
    \end{align*}
    \tcblower
    \begin{symbols}
        $\rho$ & \unit{kg/m^3} & Dichte des Bauteils \\
        $d$ & \unit{m} & Dicke des Bauteils
    \end{symbols}
\end{formelbox}

\begin{formelbox}{Doppelschalen-Resonanzfrequenz}
    \begin{align*}
        f_{res} & = \frac{60}{\sqrt{\frac{m' \cdot d}{\unit{kg/m^2} \cdot \unit{m}}}} \unit{Hz}
    \end{align*}
    \tcblower
    \begin{symbols}
        $m'$ & \unit{kg/m^2} & Reduzierte Flächenmasse \\
        $d$ & \unit{m} & Abstand zwischen den Schalen
    \end{symbols}
\end{formelbox}

\begin{formelbox}{Reduzierte Flächenmasse}
    \begin{align*}
        \frac{1}{m'} & = \frac{1}{m_1'} + \frac{1}{m_2'} \\
        m' & = \frac{m_1' \cdot m_2'}{m_1' + m_2'}
    \end{align*}
\end{formelbox}

\begin{formelbox}{Federkonstante einer Luftschicht}
    \begin{align*}
        s' & = \frac{E}{d}
    \end{align*}
    \tcblower
    \begin{symbols}
        $s'$ & \unit{N/m^3} & Federkonstante pro Fläche \\
        $E$ & \unit{Pa} & Elastizitätsmodul der Luft \\
        $d$ & \unit{m} & Dicke der Luftschicht
    \end{symbols}
\end{formelbox}

\begin{formelbox}{Doppelschalenresonanz allgemein}
    \begin{align*}
        \omega_{res} & = \sqrt{\frac{s'}{m'}} \\
        f_{res} & = \frac{\omega_{res}}{2\pi}
    \end{align*}
\end{formelbox}

\begin{formelbox}{Doppelglasscheiben mit Luftschicht}
    \begin{align*}
        f_{res} & = 1,2 \cdot \sqrt{\frac{1}{d} \cdot \left(\frac{1}{d_1} + \frac{1}{d_2}\right)} \unit{kHz}
    \end{align*}
    \begin{itemize}
        \item $\rho_{Glas} = 2500 \unit{kg/m^3}$
        \item $d$: Abstand zwischen den Scheiben [mm]
        \item $d_1, d_2$: Dicke der Scheiben [mm]
    \end{itemize}
\end{formelbox}

\begin{formelbox}{Hohlraumresonanz}
    \begin{align*}
        f_H & = \frac{n \cdot c}{2 \cdot d} \\
        d & = \frac{n \cdot \lambda}{2}
    \end{align*}
    \tcblower
    \begin{symbols}
        $n$ & --- & $n$-te Resonanz \\
        $c$ & \unit{m/s} & Schallgeschwindigkeit \\
        $d$ & \unit{m} & Dicke des Hohlraums \\
        $\lambda$ & \unit{m} & Wellenlänge
    \end{symbols}
\end{formelbox}

% ==============================================================================
% THEMA 7: LAUTHEIT UND PSYCHOAKUSTIK
% ==============================================================================

\mysection{7. Lautheit und Psychoakustik}

\begin{formelbox}{Lautheit in Sone}
    \begin{align*}
        S & = 2^{\frac{L_R - 40}{10}}
    \end{align*}
    \tcblower
    \begin{symbols}
        $S$ & \unit{sone} & Lautheit \\
        $L_R$ & \unit{phon} & Lautstärkepegel
    \end{symbols}
    \begin{itemize}
        \item Bei mehreren Geräuschen gleicher Frequenz: Gesamtpegel berechnen, dann in Sone umrechnen
        \item Bei verschiedenen Frequenzen: Einzelwerte in Phon umrechnen, dann in Sone, dann normale Addition
    \end{itemize}
\end{formelbox}

\begin{formelbox}{Äquivalenter Dauerschallpegel}
    \begin{align*}
        L_{eq} & = 10 \cdot \log\left(\frac{1}{T} \sum_{j=1}^{N} t_j \cdot 10^{0,1 \cdot L_j}\right) \\
        L_{eq} & = 10 \cdot \log\left(\sum_{j=1}^{N} q_j \cdot 10^{0,1 \cdot L_j}\right), \quad q_j = \frac{t_j}{T}
    \end{align*}
    \tcblower
    \begin{symbols}
        $L_{eq}$ & \unit{dB} & Äquivalenter Dauerschallpegel \\
        $T$ & \unit{s} & Gesamtzeitraum \\
        $t_j$ & \unit{s} & Zeitanteil des Pegels $L_j$ \\
        $q_j$ & --- & Zeitanteil
    \end{symbols}
\end{formelbox}

\begin{formelbox}{Wirkpegel}
    \begin{align*}
        L_{wirk} & = 10 \cdot \log\left(\frac{1}{T} \left((n \cdot 5\unit{s}) \cdot 10^{0,1 \cdot L_{laut}} + (T - n \cdot 5\unit{s}) \cdot 10^{0,1 \cdot L_{leise}}\right)\right)
    \end{align*}
    \tcblower
    \begin{symbols}
        $L_{wirk}$ & \unit{dB} & Wirkpegel \\
        $n$ & --- & Anzahl der lauten Peaks \\
        $5\unit{s}$ & --- & Verlängerung jedes lauten Ereignisses
    \end{symbols}
\end{formelbox}

% ==============================================================================
% THEMA 8: MECHANISCHE SCHWINGUNGEN
% ==============================================================================

\mysection{8. Mechanische Schwingungen}

\begin{formelbox}{Feder-Masse-Schwinger}
    \begin{align*}
        \omega & = \sqrt{\frac{D}{m}} \\
        f & = \frac{\omega}{2\pi} = \frac{1}{T}
    \end{align*}
    \tcblower
    \begin{symbols}
        $\omega$ & \unit{1/s} $\equiv$ \unit{Hz} & Kreisfrequenz \\
        $T$ & \unit{s} & Periodendauer \\
        $f$ & \unit{Hz} & Frequenz \\
        $D$ & \unit{N/m} & Federkonstante \\
        $m$ & \unit{kg} & Masse
    \end{symbols}
\end{formelbox}

\begin{formelbox}{Schwingungsfunktion}
    \begin{align*}
        y(t) & = y_0 \cdot \cos(\omega t + \varphi) \\
        \dot{y}(t) & = -y_0 \cdot \omega \cdot \sin(\omega t + \varphi) \\
        \ddot{y}(t) & = -y_0 \cdot \omega^2 \cdot \cos(\omega t + \varphi)
    \end{align*}
    \tcblower
    \begin{symbols}
        $y(t)$ & \unit{m} & Auslenkung \\
        $\dot{y}(t)$ & \unit{m/s} & Geschwindigkeit \\
        $\ddot{y}(t)$ & \unit{m/s^2} & Beschleunigung \\
        $y_0$ & \unit{m} & Amplitude \\
        $\varphi$ & \unit{rad} & Phasenverschiebung
    \end{symbols}
\end{formelbox}

% ==============================================================================
% THEMA 9: ELEKTROMAGNETISCHE STRAHLUNG
% ==============================================================================

\mysection{9. Elektromagnetische Strahlung}

\begin{formelbox}{Elektromagnetisches Spektrum (Licht)}
    \begin{align*}
        \text{Gammastrahlung} & < 0,005 \text{ nm} \\
        \text{Röntgenstrahlung} & 0,005 - 100 \text{ nm} \\
        \text{UV-Strahlen} & 100 - 400 \text{ nm} \\
        \text{Licht (VIS)} & 400 - 780 \text{ nm} \\
        \text{NIR} & 780 - 1400 \text{ nm} \\
        \text{MIR/FIR} & 1400 \text{ nm} - 1 \text{ mm} \\
        \text{Mikrowellen} & 1 \text{ mm} - 1 \text{ m} \\
        \text{Radio} & 1 \text{ m} - 10 \text{ km}
    \end{align*}
\end{formelbox}

\begin{formelbox}{Strahlungsintensität}
    \begin{align*}
        \vec{I} & = \vec{E} \times \vec{H} \eh{W/m^2}
    \end{align*}
    \tcblower
    \begin{symbols}
        $\vec{I}$ & \unit{W/m^2} & Intensitätsvektor \\
        $\vec{E}$ & \unit{V/m} & Elektrische Feldstärke \\
        $\vec{H}$ & \unit{A/m} & Magnetische Feldstärke
    \end{symbols}
\end{formelbox}

\begin{formelbox}{Bestrahlungsstärke}
    \begin{align*}
        B & = \frac{I}{4} \\
        \vec{B} & = \vec{I} \cdot \cos\theta
    \end{align*}
\end{formelbox}

\begin{formelbox}{Plancksche Konstanten}
    \begin{align*}
        2\pi hc^2 & = 3,741 \cdot 10^{-16} \unit{W\cdot m^2} \\
        \frac{hc}{k} & = 1,438 \cdot 10^{-2} \unit{m\cdot K}
    \end{align*}
\end{formelbox}

% ==============================================================================
% ANHANG: WICHTIGE KONSTANTEN UND WERTE
% ==============================================================================

\mysection{Anhang: Wichtige Konstanten}

\begin{formelbox}{Physikalische Konstanten}
    \begin{align*}
        \sigma & = 5,67 \cdot 10^{-8} \eh{W/(m^2 K^4)} & \text{(Stefan-Boltzmann)} \\
        R & = 8,31 \eh{J/(mol\cdot K)} & \text{(Allgemeine Gaskonstante)} \\
        R_D & = 461,5 \eh{J/(kg\cdot K)} & \text{(Wasserdampf)} \\
        c_W & = 4,182 \eh{J/(g\cdot K)} & \text{(Spezifische Wärmekapazität Wasser)}
    \end{align*}
\end{formelbox}

\begin{formelbox}{Wärmeübergangswiderstände}
    \begin{align*}
        R_{si} & = 0,10 \eh{m^2 K/W} & \text{(strömend, Wärmestrom nach oben)} \\
        R_{si} & = 0,13 \eh{m^2 K/W} & \text{(strömend, Wärmestrom horizontal)} \\
        R_{si} & = 0,17 \eh{m^2 K/W} & \text{(strömend, Wärmestrom nach unten)} \\
        R_{se} & = 0,04 \eh{m^2 K/W} & \text{(außen)}
    \end{align*}
\end{formelbox}

\begin{formelbox}{Solarkonstante}
    \begin{align*}
        q_{max} & = 1,3 \eh{kW/m^2} & \text{(oberhalb Atmosphäre)} \\
        q_{max} & = 1,0 \eh{kW/m^2} & \text{(unterhalb Atmosphäre)}
    \end{align*}
\end{formelbox}

\end{multicols*}
\end{document}